\chapter{Reference}
\label{ch:reference}

\section{Keyboard shortcuts}

\subsection{Viewport}

Single-letter shortcuts, no modifier. They yield to text fields while you
are typing.

\begin{center}
\small
\begin{tabular}{@{}ll@{}}
\toprule
Key & Action \\
\midrule
\kbd{R} & Rotation mode \\
\kbd{T} & Translation (pan) mode \\
\kbd{S} & Selection mode \\
\kbd{I} & Insertion mode \\
\kbd{D} & Distance measurement mode \\
\kbd{A} & Angle measurement mode \\
\kbd{O} & Toggle orthographic / perspective \\
\kbd{F} & Frame structure \\
\kbd{Tab} / \kbd{Shift}+\kbd{Tab} & Next / previous structure tab \\
\kbd{Delete} / \kbd{Backspace} & Delete selection (Selection mode) \\
\bottomrule
\end{tabular}
\end{center}

\begin{calnote}
Every row above except \kbd{Delete}/\kbd{Backspace} is remappable in
\menu{Edit}\menusep\menu{Preferences}\menusep\menu{Hotkeys}: a table with
one key-capture cell per action, live conflict detection against both this
list and the Application shortcuts below, and a per-row \ui{Reset} plus a
\ui{Reset All to Defaults}. \kbd{Tab}/\kbd{Shift}+\kbd{Tab} cycling wraps
around and is a no-op with a single tab open; both keys yield to the
viewport losing focus, same as the letter shortcuts above.
\end{calnote}

\subsection{Molecular Design}

Live only while the Molecular Design window (Section~\ref{sec:molecular-design})
has focus, which is what lets both it and the viewport use single letters
without either seeing the other's keys. All remappable, under
\textbf{Molecular Design} in
\menu{Edit}\menusep\menu{Preferences}\menusep\menu{Hotkeys}.

\begin{center}
\small
\begin{tabular}{@{}ll@{}}
\toprule
Key & Tool \\
\midrule
\kbd{V} & Selection \\
\kbd{1} / \kbd{2} / \kbd{3} & Single / double / triple bond \\
\kbd{4} / \kbd{5} & Wedge (bold) / hashed stereo bond \\
\kbd{C} & Chain \\
\kbd{L} & Atom label \\
\kbd{X} & Caption \\
\kbd{P} & Formal charge \\
\kbd{E} & Eraser \\
\kbd{B} & Ring template \\
\kbd{Y} & Tidy the drawing \\
\kbd{Ctrl}+\kbd{Return} & Send to 3D Viewport \\
\bottomrule
\end{tabular}
\end{center}

Undo, redo, copy, paste and select-all use the application's own bindings
inside the sketcher too, so remapping \kbd{Ctrl}+\kbd{Z} remaps it everywhere.

\subsection{Application}

\begin{center}
\small
\begin{tabular}{@{}ll@{}}
\toprule
Shortcut & Action \\
\midrule
\kbd{Ctrl}+\kbd{N} & New workspace \\
\kbd{Ctrl}+\kbd{O} & Open structure \\
\kbd{Ctrl}+\kbd{T} & Open trajectory \\
\kbd{Ctrl}+\kbd{Shift}+\kbd{O} & Open project \\
\kbd{Ctrl}+\kbd{S} & Save project \\
\kbd{Ctrl}+\kbd{Shift}+\kbd{S} & Save structure as \\
\kbd{Ctrl}+\kbd{Shift}+\kbd{T} & Save trajectory as \\
\kbd{Ctrl}+\kbd{E} & Export image \\
\kbd{Ctrl}+\kbd{W} & Close tab \\
\kbd{Ctrl}+\kbd{Q} & Quit \\
\kbd{Ctrl}+\kbd{Z} & Undo \\
\kbd{Ctrl}+\kbd{Shift}+\kbd{Z} & Redo \\
\kbd{Ctrl}+\kbd{Shift}+\kbd{A} & Add atom \\
\kbd{Ctrl}+\kbd{B} & Bond editor \\
\kbd{Ctrl}+\kbd{P} & Preferences \\
\kbd{Ctrl}+\kbd{R} & Single-point calculation \\
\bottomrule
\end{tabular}
\end{center}

On macOS, \kbd{Ctrl} in this table is \kbd{Cmd}.

\subsection{Mouse}

\begin{center}
\small
\begin{tabular}{@{}ll@{}}
\toprule
Gesture & Action \\
\midrule
Left drag                  & Depends on interaction mode \\
Middle drag                & Pan (any mode) \\
\kbd{Shift}+left drag      & Pan (any mode) \\
Wheel                      & Zoom \\
Left click                 & Pick atom \\
\kbd{Ctrl}/\kbd{Cmd}+click & Toggle atom in selection \\
Double click               & Frame structure \\
\bottomrule
\end{tabular}
\end{center}

\section{Menu map}

\begin{description}
  \item[\menu{File}] New Workspace $\cdot$ Open (Structure, Trajectory)
    $\cdot$ Save (Structure As, Trajectory As) $\cdot$ Import / Export
    (Image, Animation, Ray-Traced Render) $\cdot$ Project Workspace (Open,
    Save, Save As) $\cdot$ Close Tab $\cdot$ Quit
  \item[\menu{Edit}] Undo $\cdot$ Redo $\cdot$ Add Atom $\cdot$ Selection
    (Change Element, Translate) $\cdot$ Delete Selected Atoms $\cdot$ Bond
    Editor $\cdot$ Unit Cell (Create Supercell) $\cdot$ Preferences
  \item[\menu{View}] Orthographic $\cdot$ Overlays (Show Unit Cell) $\cdot$
    Alignment (Frame Structure, XY, XZ, YZ) $\cdot$ Visual Effects $\cdot$
    panel toggles
  \item[\menu{Build}] Nanomaterial Builder $\cdot$ Surface Slab $\cdot$
    Metallic Nanoparticle (Wulff) $\cdot$ Special Quasirandom Structure
    (SQS) $\cdot$ Normal Modes / Phonon Builder $\cdot$ From Database
    $\cdot$ Structure Perturbation / Noise
  \item[\menu{Simulation}] New Calculation (ASE) $\cdot$ New Remote
    Calculation $\cdot$ Electronic Bands / PDOS $\cdot$ Dataset Manager
    (MLIP)
  \item[\menu{Analysis}] Structure Factor S(q) $\cdot$ X-Ray Diffraction
    $\cdot$ Radial Distribution Function $\cdot$ Bond Length / Angle
    Distributions $\cdot$ Coordination Numbers (CN / GCN) $\cdot$
    Warren-Cowley analysis $\cdot$ Local Entropy Analysis $\cdot$ Raman
    Modes $\cdot$ Volumetric Data $\cdot$ Adsorption \& Catalysis $\cdot$
    Brillouin Zone / k-Path
  \item[\menu{Help}] Documentation $\cdot$ GitHub Repository $\cdot$ About
    Calango
\end{description}

\section{Python dependencies by feature}

\begin{center}
\small
\begin{tabular}{@{}ll@{}}
\toprule
Package & Needed for \\
\midrule
\texttt{ase}, \texttt{numpy} & Everything (structure I/O, building, jobs) \\
\texttt{spglib}       & Symmetry detection, Raman modes \\
\texttt{paramiko}     & HPC panel \\
\texttt{pillow}       & GIF animation export \\
\texttt{imageio}, \texttt{imageio-ffmpeg} & MP4 animation export \\
\texttt{mace-torch}   & MACE machine-learning potentials \\
\texttt{gpaw}         & DFT band structures and PDOS \\
\texttt{icet}         & Preferred SQS backend (optional) \\
\bottomrule
\end{tabular}
\end{center}

External binaries: \texttt{pw.x} (Quantum ESPRESSO), \texttt{siesta},
VASP, \texttt{orca}, \texttt{povray} or \texttt{tachyon}.

\section{Job directory contents}

A typical job directory under \file{.calango\_tmp/}
(Section~\ref{sec:calango-tmp}):

\begin{center}
\small
\begin{tabular}{@{}ll@{}}
\toprule
File & Written by \\
\midrule
\file{run.py}            & The generated (possibly hand-edited) script \\
\file{structure.extxyz}  & The staged input geometry \\
\file{opt.traj}          & Geometry optimization trajectory \\
\file{md.traj}           & Molecular dynamics trajectory \\
\file{optimized.extxyz}  & Final relaxed geometry \\
\file{bands.json}, \file{pdos.json} & Electronic structure results \\
\file{phonon\_bands.csv}, \file{phonon\_dos.csv} & Phonon results \\
\file{perturbed.extxyz}  & Noise ensemble checkpoint \\
\file{job.sh}            & Scheduler wrapper (remote jobs) \\
\file{calango\_job.out}, \file{calango\_job.err} & Remote stdout/stderr \\
\bottomrule
\end{tabular}
\end{center}

\section{Progress markers}

Generated scripts communicate with the interface by printing markers to
standard output. If you hand-edit a script or write your own, emitting
these keeps the Job panel live:

\begin{center}
\small
\begin{tabular}{@{}ll@{}}
\toprule
Marker & Effect \\
\midrule
\texttt{CALANGO\_PROGRESS step total}   & Progress bar \\
\texttt{CALANGO\_ENERGY step value}     & Energy plot (eV) \\
\texttt{CALANGO\_TEMP step value}       & Temperature plot (K) \\
\texttt{CALANGO\_FMAX step value}       & Force plot (eV/\AA) \\
\texttt{CALANGO\_PRESSURE step value}   & Pressure plot (GPa) \\
\texttt{CALANGO\_TARGET\_TEMP value}    & Thermostat reference line \\
\texttt{CALANGO\_TARGET\_PRESSURE value}& Barostat reference line \\
\texttt{CALANGO\_CELL} + \texttt{CALANGO\_FRAME} & Live trajectory streaming \\
\texttt{CALANGO\_RESULT \ldots}         & Final summary line \\
\bottomrule
\end{tabular}
\end{center}

\section{Troubleshooting}

\begin{description}
  \item[Structures will not open; job features are disabled.]
    ASE is not importable. Run \lstinline|calango --probe-python| to see
    which interpreter is being used, then either install ASE there or set
    \opt{CALANGO\_PYTHON} to an environment that has it
    (Section~\ref{sec:python-env}).

  \item[The space group is \texttt{P1} for an obviously symmetric crystal.]
    Numerical noise in the coordinates. Raise \ui{Tolerance} in the
    \ui{Structure} panel --- a relaxed cell often needs 0.01\,\AA{} or
    more.

  \item[Bonds are missing or spurious.]
    Adjust \ui{Covalent cutoff multiplier} in the bond editor
    (Section~\ref{sec:bonds}), or add and suppress individual bonds by
    hand. Metallic and ionic systems rarely match a covalent-radius rule.

  \item[Double and triple bonds do not appear automatically.]
    By design. Assign them with the \ui{Bond order} buttons
    (Section~\ref{sec:bond-orders}).

  \item[A job fails immediately with an import error.]
    The job is running in a different environment from Calango itself.
    Check the \ui{Execution Environment} status line in the calculation
    dialog (Section~\ref{sec:exec-env}).

  \item[The GPAW backend is greyed out.]
    GPAW is not importable in the selected environment. Point the
    \ui{Environment} selector at a Conda environment where it is installed.

  \item[The remote panel reports a connection failure.]
    Test the same host, port, user and key with \texttt{ssh} from a
    terminal first. Remember that with a \ui{Key file} set, the
    \ui{Password} field doubles as that key's passphrase rather than a
    login password (Chapter~\ref{ch:remote} explains the full inference).

  \item[Adsorption site detection finds no bridge or hollow sites.]
    The surface cell is too small to have in-plane neighbours. Build a
    supercell of the slab first.

  \item[The viewport looks soft after enabling depth of field.]
    Expected --- the effect trades multisample anti-aliasing for the blur
    pass (Section~\ref{sec:visual-effects}). Disable it before exporting
    still images.

  \item[Animation export fails.]
    Install \texttt{pillow} for GIF, or \texttt{imageio} and
    \texttt{imageio-ffmpeg} for MP4, in Calango's Python environment.
\end{description}

\section{Citing Calango}
\label{sec:citing}

\menu{Help}\menusep\menu{About Calango} has a \ui{Citations} tab, right
after \ui{Third-Party Licenses}: it groups the references worth citing for
a Calango-produced result into \ui{Core} (Calango itself, and ASE, which
every generated script builds on), \ui{Databases} (Materials Project,
PubChem, C2DB --- whichever a structure actually came from), \ui{Calculators}
(the DFT/MD engine a wizard generated a script for --- cite the ones a
given result actually ran on, not the whole list), and \ui{Libraries}
(spglib, NumPy, phonopy, icet, dftd4 --- customary to cite where a venue
expects it). Each entry has a \ui{BibTeX viewer\ldots} button opening a
read-only, monospaced, copyable BibTeX record, so citing a result correctly
never means retyping a reference by hand or guessing a citation key.

\section{Further reading}

\begin{description}
  \item[\file{README.md}] Feature overview and build instructions.
  \item[\file{docs/sphinx/}] The Sphinx manual --- the primary, most
    frequently updated reference, covering every builder, wizard,
    calculator and analysis viewer in more depth than this guide.
  \item[\file{docs/packaging.pdf}] Building the macOS and Linux installers.
  \item[\url{https://github.com/seixas-research/calango}] Source, issues
    and releases.
  \item[\url{https://wiki.fysik.dtu.dk/ase/}] ASE documentation --- the
    reference for anything in a generated script.
\end{description}
